\chapter{What leave-rule-out does not fix}
\label{ch:limits}

\begin{saybox}[title={What to say at the start of this section}]
Volunteer these limitations before the panel raises them. The credibility
of everything in the previous chapters depends on being seen to understand
what the revised protocol still cannot measure. There are four items, and
the fourth --- censoring --- is the one that was found during this audit
and was not previously documented.
\end{saybox}

\section{Limitation 1: rules are not independent of each other}

Leave-rule-out removes one \texttt{source\_rule\_id} at a time. It does not
remove the \emph{components} that rule shares with others.
\Cref{ch:rules} showed that rule identifiers are pipe-delimited
composites: \texttt{NATAMYCIN\_REVIEW} appears in three different soft
rules, \texttt{HARD\_SEMI\_REVIEW} appears in every hard/general rule
without exception, and \texttt{CLOSTRIDIA\_GRANA} appears in every
clostridia rule in two specialists.

When \texttt{RICOTTA\_MAP|NATAMYCIN\_REVIEW} is held out, the training set
still contains \texttt{NATAMYCIN\_REVIEW} alone and
\texttt{HMMC\_NAT|NATAMYCIN\_REVIEW}. If those rules share calibration
constants or a common parametric form --- which their shared naming makes
plausible but which cannot be verified without the generator --- then the
fold is not a clean test of transfer to an unseen mechanism.

\begin{limitbox}[title={Consequence}]
Leave-rule-out is best understood as an \emph{upper bound relaxation}, not
an exact measurement of mechanism transfer. It is strictly more honest than
the context-holdout protocol, and it is still optimistic relative to a
test on a genuinely independent mechanism.

A stricter variant --- leave-one-\emph{source-component}-out, withholding
every rule containing a given component --- is implementable from the
identifiers alone and is recommended in \Cref{ch:proposals}. It has not
been run.
\end{limitbox}

\section{Limitation 2: all rules come from the same generator}

Even a perfectly independent set of rules would share the generator's
global assumptions: its functional forms, its noise model, its treatment of
temperature dependence, its clipping and bounding behaviour. A model that
transfers across rules from one generator has demonstrated that it can
learn that generator's family of relationships. It has not demonstrated
anything about relationships the generator cannot produce.

Since the generator is not in this repository (\Cref{sec:limit-evidence}),
the size of this gap cannot be quantified here at all. It can only be
flagged.

\section{Limitation 3: synthetic transfer is not real-world transfer}

This is the most important non-claim in the report, and it should be stated
in exactly these terms at the meeting:

\begin{findingbox}[title={The claim that may be made, and the one that may not}]
\textbf{May be claimed:} under leave-one-generation-rule-out
cross-validation on the V7 synthetic corpus, LightGBM achieves a median
fold $\rsq$ of $0.780$ and a pooled out-of-fold $\rsq$ between $0.67$ and
$0.96$ depending on the specialist, with median absolute error under four
days.

\textbf{May not be claimed:} that the model predicts the shelf life of
real cheese to that accuracy. No internal metric computed on synthetic
data can support that statement, however the folds are constructed. The
only evidence that bears on it is external, and \Cref{ch:external}
presents it.
\end{findingbox}

\section{Limitation 4: right-censoring in the safety endpoints}
\label{sec:censoring}

This was discovered during the present audit and is, on the evidence, not
documented elsewhere in the project.

The column \texttt{target\_is\_lower\_bound} flags rows whose recorded
\texttt{shelf\_life\_days} is a \emph{lower bound} on the true event time
rather than the event time itself --- the standard situation in a challenge
study where the pathogen had not reached its criterion when the study
ended. Its distribution across the six specialists:

\begin{center}\small
\begin{tabular}{lrrr}
\toprule
Specialist & rows & censored ($=1$) & share censored \\
\midrule
soft / general & 8\,800 & 0 & $0\%$ \\
semi-hard / general & 7\,200 & 0 & $0\%$ \\
hard / general & 6\,000 & 0 & $0\%$ \\
\addlinespace[2pt]
soft / safety & 4\,800 & 3\,469 & $72.3\%$ \\
semi-hard / safety & 4\,000 & 3\,784 & $94.6\%$ \\
hard / safety & 3\,200 & 3\,200 & $\mathbf{100.0\%}$ \\
\bottomrule
\end{tabular}
\end{center}

\begin{findingbox}[title={Every row of the hard/safety-endpoint specialist is right-censored}]
All $3\,200$ rows of the hard safety specialist carry
\texttt{target\_is\_lower\_bound = 1}. The specialist is therefore not
trained to predict the time at which a pathogen criterion is reached. It is
trained to predict \emph{the last time point at which the criterion had not
yet been reached} --- a quantity determined partly by the duration of the
underlying study design.

For soft/safety the two subpopulations differ sharply: censored rows
average $39.3$ days, uncensored rows $12.1$ days. A squared-error
regression fitted across both is fitting a mixture of two different
quantities.
\end{findingbox}

\subsection{Why this matters statistically}

For a censored observation the recorded value satisfies
$y_{\text{obs}} = \min(T, C)$ where $T$ is the true event time and $C$ the
censoring time. Least-squares regression minimises
$\sum_i (y_{\text{obs},i} - \hat{f}(\mathbf{x}_i))^2$, which targets
$\mathbb{E}[\min(T, C) \mid \mathbf{x}]$, not $\mathbb{E}[T \mid \mathbf{x}]$.
Since $\min(T,C) \le T$ always, the fitted model is biased low as an
estimator of the true failure time --- and the magnitude of that bias
depends on the censoring distribution, which is an artefact of study design
rather than of cheese microbiology.

Recall also from \Cref{sec:exclusions} that \texttt{target\_is\_lower\_bound}
is on the exclusion list, for the defensible reason that it would not be
known at prediction time. The consequence is that the loss function has no
way to distinguish the two kinds of row.

\subsection{What this does not mean}

It does not mean the safety predictions are worthless. A model of
``time by which the criterion was demonstrably not yet reached'' is a
conservative quantity, and for a safety application conservatism in this
direction is arguably the safer error. But it must be \emph{described}
correctly, and at present it is described as shelf life.

Standard remedies exist --- Tobit regression, accelerated failure time
models, or a censoring-aware gradient-boosting objective --- and
\Cref{ch:proposals} lists the diagnostic that should be run first.

\section{Limitation 5: hyperparameter selection was not group-aware}
\label{sec:hparam}

Under the original protocol, LightGBM and XGBoost select their number of
boosting rounds by early stopping on the validation partition. That
partition is drawn from the same generation rules as the test partition
(\Cref{ch:discovery}, $100\%$ rule overlap). The stopping iteration is
therefore selected using information about the rules on which the model is
subsequently scored.

The effect is second-order compared with the partitioning issue itself ---
early stopping selects one scalar, not a model --- but it is a real
dependency and it is worth stating. In the leave-rule-out harness the
rounds are fixed a priori, which removes the dependency but leaves $M$
untuned. A fully clean design would use nested, group-aware selection:
an inner leave-rule-out loop over the training rules of each outer fold.
This has not been run.

\section{Limitation 6: feature audit}
\label{sec:leakage}

Every feature reaching the model was checked against the target for the
classical leakage patterns. The findings:

\begin{itemize}[leftmargin=1.3em, itemsep=2pt]
\item \textbf{No identifier reaches the model.} \texttt{row\_id},
  \texttt{context\_id}, \texttt{formulation\_id} and
  \texttt{source\_rule\_id} are all excluded.
\item \textbf{No provenance column reaches the model.}
  \texttt{data\_origin}, \texttt{source\_url}, \texttt{source\_note},
  \texttt{evidence\_class}, \texttt{generation\_rule\_version} are
  excluded.
\item \textbf{No recomputed target reaches the model.}
  \texttt{shelf\_life\_days\_v4\_original} and
  \texttt{target\_recalibration\_factor} --- which would be direct
  leakage --- are on the exclusion list.
\item \textbf{\texttt{indicator\_threshold} and
  \texttt{initial\_indicator\_value} are included, and should be.} They are
  legitimate inputs: the threshold defines what ``end of shelf life'' means
  for that row, and the initial load is a measurable property of the
  product at time zero. Both would be available at prediction time.
\item \textbf{\texttt{matrix\_ripening\_days} is included.} It is a
  composition descriptor, not a downstream measurement, but for hard
  cheeses it correlates with the maturation state that also drives the
  target. It is not leakage, but it is the feature most worth a second
  look.
\end{itemize}

No target-derived feature was found among the model inputs. The high scores
are not explained by leakage in the technical sense.
